\documentclass[10pt,a4paper]{amsart}
\usepackage[margin=19mm]{geometry}
\usepackage{amsmath,amssymb,mathtools}
\usepackage[scr=rsfs,cal=euler]{mathalfa}
\usepackage{xfrac}
\usepackage[unicode,hidelinks,bookmarks=false]{hyperref}
\newcommand{\ie}{i.e.\ }
\newcommand{\cf}{cf.\ }
\newcommand{\resp}{respectively\ }
\newcommand{\Subsection}{\subsection}
\providecommand{\nobreakdash}{\nobreak\hbox{-}}
\setlength{\emergencystretch}{2em}
\multlinegap0pt
\usepackage{amssymb}
\usepackage{algorithm}\usepackage{algpseudocode}
\newtheorem{lemma}{Lemma}
\newcommand{\C}{{\mathbb C}}
\newcommand{\E}{{\mathbb E}}
\newcommand{\tr}{{\rm tr}}
\newcommand{\ve}{{\mathbf e}}
\newcommand{\vx}{{\mathbf x}}
\title{The Fusion Frame Phase Retrieval}
\author{Haixia Liu, Bing Gao, Yang Wang}
\date{Source: arXiv:2609.08531v1 (CC BY 4.0)}
\begin{document}
\maketitle

\noindent\emph{Source and licence.} The Fusion Frame Phase Retrieval. Version arXiv:2609.08531v1 (CC BY 4.0), distributed by arXiv under the Creative Commons Attribution 4.0 International licence, http://creativecommons.org/licenses/by/4.0/, as recorded in the arXiv licence metadata for this exact version and shown on the abstract page https://arxiv.org/abs/2609.08531v1. Full licence notice: \texttt{licenses/arxiv-2609.08531-CC-BY-4.0.txt}.\par\smallskip

\noindent\textit{Editorial scope.} Counted unit: the article's lemma lem:expectations (source0.tex lines 960--976). The article's complete original proof of it is delivered with it (source0.tex lines 978--1073, one \texttt{proof} environment of 96 lines). No mathematics is written, repaired or re-derived: the statement and the proof are the article's own lines, in the article's own notation, and no retained line is substituted or re-flowed.  Retained uncounted as the source-local context the proof needs: lines 873--887; lines 924--940.  Nothing was omitted from any retained span, so the package carries no omission record.  No retained line contains a citation, so the wrapper carries no bibliography and no bibliography entry.  No retained line contains a figure environment, an image-inclusion command or drawing code, so no image is delivered and the package declares no asset dependency.  Editorial formatting only, in addition: an AMS \texttt{amsart} preamble that restates the article's own declarations and macros used by the retained lines, namely the environments \texttt{lemma} on separate counters, together with the standard packages those lines need.  The retained environments print in this document as Lemma 1 in order of appearance, whereas the article prints the same items with the numbering of its own full document.  The complete native source of the article as distributed in the arXiv e-print for this exact version is bundled unchanged as \texttt{sources/209722/source0.tex}.
	\begin{lemma}[Scalar distribution and moments]
		\label{lem:beta-d}
		Let $A$ be an orthogonal projection in $\C^d$ distributed from the Haar measure on the set of rank-$r$ orthogonal projection matrices, where $1\le r<d$. Then, for every fixed nonzero vector $\vx\in\C^d$,
		\[
		\frac{\vx^*A\vx}{\|\vx\|^2}\sim\operatorname{Beta}(r,d-r).
		\]
		Consequently, for every integer $p\ge 1$,
		\[
		\E\left[(\vx^*A\vx)^p\right]=\|\vx\|^{2p}\frac{(r)_p}{(d)_p}\leq\|\vx\|^{2p}\Gamma(p+1)\left(\frac{r}{d}\right)^p,
		\]
		where
		\[
		(r)_p=r(r+1)\cdots(r+p-1),\qquad (d)_p=d(d+1)\cdots(d+p-1).
		\]
	\end{lemma}

	\begin{lemma}
		\label{lem:expectation1}
		Let $A$ be a rank-$r$ orthogonal projection in $\C^d$ distributed from the Haar measure, where $1\le r<d$. Then:
		\begin{enumerate}
			\item For any $i\ne j$,
			\[
			\E(a_{ij}^2)=\E(a_{12}^2),\qquad \E(|a_{ij}|^2)=\E(|a_{12}|^2).
			\]
			\item For every unitary matrix $P\in\C^{d\times d}$ and every $\vx\in\C^d$, define
			$
			\mathcal E(\vx)=\E_A\big[(\vx^*A\vx)A\big]$. 
			If $\widetilde{\vx}=P\vx$, then
			\[
			\mathcal E(\vx)=P^*\mathcal E(\widetilde{\vx})P.
			\]
		\end{enumerate}
	\end{lemma}

	\begin{lemma}
		\label{lem:expectations}
		Let $A$ be a rank-$r$ orthogonal projection in $\C^d$ distributed from the Haar measure, where $1\le r<d$. Then
		\[
		\E A=\frac rd I_d.
		\]
		Moreover, for every fixed $\vx\in\C^d$,
		\[
		\E\big[(\vx^*A\vx)A\big]=\mu\|\vx\|^2I_d+\zeta\vx\vx^*,\qquad\E\big[A\vx(A\vx)^*\big]=\zeta\|\vx\|^2I_d+\mu\vx\vx^*,\qquad\E\big[A\vx(A\vx)^\top\big]=\lambda\vx\vx^\top,
		\]
		where
		\begin{equation}\label{defi_para}
			\lambda=\frac{r(r+1)}{d(d+1)},\qquad
			\mu=\frac{dr^2-r}{d(d+1)(d-1)},\qquad
			\zeta=\lambda-\mu=\frac{r(d-r)}{d(d+1)(d-1)}.
		\end{equation}
	\end{lemma}

	\begin{proof}
		For every unitary matrix $P$, the random matrices $A$ and $PAP^*$ have the same distribution. Hence
		\[
		P(\E A)P^*=\E A.
		\]
		Thus $\E A$ commutes with every unitary matrix and must be a scalar multiple of the identity. Since $\tr(A)=r$,
		\[
		\E A=\frac rd I_d.
		\]
		
		We first consider $\vx=\ve_1$. Set
		\[
		M_1=\E\big[(\ve_1^*A\ve_1)A\big].
		\]
		For every unitary matrix $Q\in\C^{(d-1)\times(d-1)}$, let $P=\operatorname{diag}(1,Q)$. Since $P\ve_1=\ve_1$, Lemma~\ref{lem:expectation1} gives $M_1=P^*M_1P$. It follows that
		\[
		M_1=\begin{bmatrix}\lambda&0\\0&\mu I_{d-1}\end{bmatrix}.
		\]
		By Lemma~\ref{lem:beta-d},
		\[
		\lambda=\E(a_{11}^2)=\frac{r(r+1)}{d(d+1)}.
		\]
		Moreover,
		\[
		\tr(M_1)=\E\big[a_{11}\tr(A)\big]=r\E(a_{11})=\frac{r^2}{d}.
		\]
		Therefore,
		\[
		\lambda+(d-1)\mu=\frac{r^2}{d},
		\]
		which yields
		\[
		\mu=\frac{dr^2-r}{d(d+1)(d-1)}.
		\]
		Since $\zeta=\lambda-\mu$, we obtain
		\[
		M_1=\mu I_d+\zeta \ve_1\ve_1^*.
		\]
		
		For a general nonzero vector $\vx$, choose a unitary matrix $P$ such that $P\vx=\|\vx\|\ve_1$. By Lemma~\ref{lem:expectation1} and the homogeneity of $\mathcal E$,
		\[
		\E\big[(\vx^*A\vx)A\big]=P^*\mathcal E(\|\vx\|\ve_1)P
		=\mu\|\vx\|^2I_d+\zeta\vx\vx^*.
		\]
		The identity is trivial for $\vx=0$.
		
		Next, set
		\[
		M_2=\E\big[A\ve_1(A\ve_1)^*\big].
		\]
		The same invariance argument gives
		\[
		M_2=\begin{bmatrix}\lambda&0\\0&\zeta I_{d-1}\end{bmatrix}.
		\]
		Indeed, its $(1,1)$-th entry is $\E(a_{11}^2)=\lambda$, while
		\[
		\tr(M_2)=\E\|A\ve_1\|^2=\E(\ve_1^*A^2\ve_1)=\E(a_{11})=\frac rd.
		\]
		Thus
		\[
		\lambda+(d-1)\zeta=\frac rd,
		\]
		and consequently
		\[
		\zeta=\frac{r(d-r)}{d(d+1)(d-1)}.
		\]
		Since $\lambda-\zeta=\mu$, we have
		\[
		M_2=\zeta I_d+\mu \ve_1\ve_1^*.
		\]
		Applying unitary invariance as above yields
		\[
		\E\big[A\vx(A\vx)^*\big]=\zeta\|\vx\|^2I_d+\mu\vx\vx^*.
		\]
		
		Finally, let
		\[
		M_3=\E\big[A\ve_1(A\ve_1)^\top\big].
		\]
		For a diagonal unitary matrix $D=\operatorname{diag}(1,e^{i\theta_2},\ldots,e^{i\theta_d})$, the matrices $A$ and $DAD^*$ have the same distribution. Hence
		\[
		M_3=DM_3D^\top.
		\]
		Since the phases $\theta_2,\ldots,\theta_d$ are arbitrary, all entries of $M_3$ vanish except possibly its $(1,1)$-th entry. As
		\[
		(M_3)_{11}=\E(a_{11}^2)=\lambda,
		\]
		we obtain
		\[
		M_3=\lambda \ve_1\ve_1^\top.
		\]
		A final application of unitary invariance gives
		\[
		\E\big[A\vx(A\vx)^\top\big]=\lambda\vx\vx^\top.
		\]
	\end{proof}

\end{document}
